\section{The 1.1.0 GEMM mechanisms}
\label{sec:mechanisms}

Section \ref{sec:gap} ends on a hypothesis: the residual gap to cuBLAS is mostly
kernel selection. The four mechanisms in this section are what test it. All four
are implemented, tested and profiled at compile time; none of them has a
throughput number yet, because the tile sweep that produces those numbers runs
under the locked clock protocol and that is an owner action. Every performance
column below therefore reads \emph{pending}, and the generator that writes these
tables leaves it that way rather than filling it from a model. The mechanisms are
documented in full in \texttt{docs/gemm.md}.

\subsection{The tile family}

Rung 6 is one block shape, 128 by 128 with a K step of 32. The family templates
the mainloop on the tuple of block extents and the warp arrangement, and ships six
instantiations. Every warp tile is 64 by 32, which is why the accumulator stays at
64 registers across the whole family: the shapes differ in shared memory footprint
and in how badly they quantize against 48 SMs, not in register pressure. Nothing
spills. The register and shared memory columns in Table \ref{tab:tiles} come from
\texttt{nvcc -{}-resource-usage} at sm\_120, so they are compile time facts and are
filled in; the throughput column is the one waiting on the sweep. The shape
constraints are \texttt{static\_assert}s rather than comments: three stages inside
the 99 KB per block opt-in, warps inside the 48 an SM carries, the staging work
dividing evenly across the block, and a warp tile that the \texttt{m16n8k16} shape
divides.

Shipping a family rather than a constant is the Nath, Tomov and Dongarra
argument~\cite{nath2010}: which tile wins is a property of the shape and the part,
and the way to find out is to instantiate them all and measure.

\begin{table}[htbp]
\centering
\begin{tabular}{llrrrr}
\toprule
tile & warps & threads & shared bytes, 3 stages & GFLOP/s & percent of cuBLAS \\
\midrule
128x128x32 & 2 by 4 & 256 & 49,152 & pending & pending \\
128x256x32 & 2 by 8 & 512 & 73,728 & pending & pending \\
256x128x32 & 4 by 4 & 512 & 73,728 & pending & pending \\
128x64x64 & 2 by 2 & 128 & 73,728 & pending & pending \\
64x128x64 & 1 by 4 & 128 & 73,728 & pending & pending \\
128x128x64 & 2 by 4 & 256 & 98,304 & pending & pending \\
\bottomrule
\end{tabular}

\caption{The tile family. Registers and shared memory are compile time facts from
the resource usage report; relative throughput awaits the locked clock tile sweep.}
\label{tab:tiles}
\end{table}

\subsection{Predicated tails}

Before 1.1.0 a shape not divisible by the tile rerouted to the WMMA kernel, and
one not divisible by 64 by 64 by 16 rerouted again to a scalar kernel with one
thread per output element. Nothing warned, and no benchmark measured what that
cost. The family predicates its edges instead, so every shape runs the same
mainloop with masked edges. The staging uses \texttt{cp.async}'s source size
field: a sixteen byte chunk that is partly past an edge copies the bytes that
exist and zero fills the rest, and one wholly past an edge copies nothing and zero
fills all sixteen. Zeros in the staged tile contribute nothing to the dot product,
so the mainloop carries no edge logic at all.

One case defeats \texttt{cp.async} outright, and it is worth naming because it is
not obvious. A sixteen byte copy needs a sixteen byte aligned source, and a row of
A starts at a multiple of $k$, so an odd $k$ puts every second row on an odd
address. Those shapes stage through plain loads and shared stores, chosen by a
flag the launcher computes from the leading dimensions and the operand base
addresses. The pipeline structure does not change: the barrier that already
separates the stage writing a buffer from the iteration reading it covers a
synchronous store just as well as an asynchronous one. The scalar half precision
path still exists but \texttt{Algo::kAuto} never routes to it, and \texttt{chosen}
reports whatever did run, so there is no silent fallback. Cost of the predicated
tails relative to the aligned path: pending.

\subsection{Split-K}

A 1024 by 1024 output at a 128 by 128 tile is 64 blocks on a 48 SM part, which is
the 0.67 wave row of Table \ref{tab:gap}: most of the machine idles for most of
the call. Split-K cuts the contraction into slices, runs each as its own grid
layer, and sums the slices in a fixup pass.

Two passes rather than atomics into C, for two reasons. Atomics would make the
result depend on the order blocks happened to retire, which is not something a
correctness test can pin down. And they would have to read C to apply beta from
inside the mainloop, which breaks the rule that beta equal to zero does not read
C. So the first pass writes raw partial sums into its own plane and never touches
C, and the second sums the planes, multiplies by alpha once, and adds beta times C
once.

The scratch is $\text{slices} \times m \times n$ floats.
\texttt{ckl::gemm\_workspace\_size} answers with it, and a Context that already
owns a workspace that large keeps the allocation out of the call; without one the
driver allocates and frees its own around the launch, which is correct but is an
allocation a benchmark must not time. The slice count a caller asks for is a
request, because each slice is rounded up to a whole number of the tile's K step,
and \texttt{ckl::gemm\_split\_k\_slices} reports how many were actually used.
Speedup at the shapes where waves are scarce: pending.

\subsection{Stream-K}

Split-K cuts every tile the same way and pays a full extra pass over C for it.
Stream-K only touches the tiles that would have left SMs idle~\cite{osama2023}.
The grid is a block of persistent CTAs followed by ordinary data parallel ones:
the data parallel CTAs take the tiles that fill whole waves, one each, and the
stream-K CTAs share out the K iterations of what is left, which is at most two
waves of tiles.

A CTA whose share stops before the end of a tile holds only part of the sum, so it
publishes its partial into its own plane of the workspace, fences, and raises a
flag. The CTA that covers the tile's last K iteration waits on the flag of every
peer whose share ended inside that tile, adds the planes into its accumulator, and
writes C once.

The direction of that wait is the safety argument and it is deliberate: an owner
only ever waits on CTAs with a lower index than its own. Blocks are dispatched in
increasing index order, so a resident owner implies every peer it can wait on is
already resident or already retired. Waiting upward would require every stream-K
CTA to be co-resident, which is an assumption about occupancy that a profiler's
instrumentation or a future register allocation change could quietly turn into a
hang. The data parallel CTAs never wait at all.

Three orderings hold the protocol together and all three are explicit. A
\texttt{\_\_threadfence} between the partial stores and the flag makes the partial
visible device wide before the flag that advertises it. A
\texttt{\_\_syncthreads} between them makes that true for every thread of the
block and not just the one raising the flag. And the persistent loop carries a
\texttt{\_\_syncthreads} of its own, because the shared tile buffers are reused
from one tile to the next and the next tile's \texttt{cp.async} stores would
otherwise race the previous tile's \texttt{ldmatrix} reads. That last one is not a
guess: it is the barrier \texttt{compute-sanitizer -{}-tool racecheck} reports when
it is removed. Speedup on small and odd shapes: pending.

\subsection{Dispatch}

\texttt{Algo::kAuto} on an FP16 input with an FP32 output now scores every tile in
the family against the shape, using wave quantization and the shared memory
budget, and selects a tile, a decomposition, and a slice count. The heuristic
reads the tile sweep CSV when one is present. The committed sweep carries
\texttt{clock\_locked=false}, which means the heuristic must not be tuned from it,
and it is not. Table \ref{tab:decomp} is the decomposition summary; like the tile
table, its performance columns wait on the sweep.

\begin{table}[htbp]
\centering
\begin{tabular}{lllr}
\toprule
decomposition & shapes it targets & what it fixes & measured speedup \\
\midrule
data parallel & one CTA per output tile & the baseline decomposition & measured (the ladder rows above) \\
predicated tails & shapes the tile does not divide & one mainloop for every shape, no silent reroute & pending \\
split-K & few waves, long contraction & idle SMs when the output is smaller than one wave & pending \\
stream-K & few waves, any contraction & the same, without a full extra pass over C & pending \\
\bottomrule
\end{tabular}

\caption{The decomposition mechanisms and what each one is expected to fix.
Measured speedups are pending the locked clock sweep; the expectations are stated
as hypotheses and are not evidence.}
\label{tab:decomp}
\end{table}
